\section{Why the proof is exact}

This chapter explains why the proof uses exact identities rather than
estimates, probability, or transcendence statements.

\subsection{The exact ledger}

The core identity is
\begin{equation*}
  2^{\wordWeight(w)}\,\wordOrbit(w,x_0)
    =5^{\operatorname{len}(w)}\,x_0+\wordA(w).
\end{equation*}
Every later equation in the proof is a projection of this identity:
block equations, segment equations, reset equations, and the CRT
equation.  There is no place where a size estimate replaces an
equality.

\subsection{Where loose estimates would fail}

An average-growth estimate would conclude that the orbit should
increase on average, which is compatible with both unboundedness and
with a long excursion followed by a cycle.  The proof needs the
opposite: a contradiction that is sensitive to a single unit of
2-adic valuation.  The decisive-window thresholds are separated by
exactly one:
\begin{equation*}
  2j+10<2j+11<2j+12.
\end{equation*}
An estimate that is off by one cannot distinguish the reachable upper
bound from the block-layer lower bound.

\subsection{The counterexample as a failure mode}

The arithmetic witness
\begin{equation*}
  s=2^{83}-3\cdot5^{35}
\end{equation*}
satisfies every size, parity, and coprimality condition used by the
old bound, but its valuation is $83$, one more than $82$.  The witness
shows exactly what an arithmetic-only argument can miss: the
reachability fields.  No estimate can repair this, because the failure
is an exact valuation mismatch.

\subsection{What probability cannot provide}

A probabilistic model can predict the distribution of step weights,
but the proof requires a statement about the specific orbit of $7$.
Equidistribution would not rule out a cycle with an exceptional
valuation profile; the proof must rule out every cycle by exact
arithmetic.

\subsection{What transcendence cannot provide}

The margin $\log_2 5$ is irrational, but the proof never needs
irrationality or transcendence.  The PMI identity is cleared of
denominators:
\begin{equation*}
  \sum_{j=0}^{L-1}2^{\wordWeight_j}\,5^{L-j}=5\,\wordA(w).
\end{equation*}
All inequalities are between integers or between rationals with
explicit denominators.  Transcendence estimates would be both too weak
and unnecessary.

\subsection{The role of the finite base}

The only finite data are the first 17 states of the orbit of $7$.  The
finite base is used to identify the first big step; every later claim
is exact arithmetic.  The base is expanded explicitly in Lean, so no
search or code generation is trusted.

\subsection{Exactness checklist}

\begin{itemize}[leftmargin=*]
\item every block equation is the word-orbit identity;
\item every segment exclusion uses exact equations and exact
  congruences;
\item the CRT candidate is the least solution of an exact linear
  equation;
\item the size conditions are exact inequalities;
\item the window thresholds are exact and never weakened;
\item the pending Lean statement is the only unverified node.
\end{itemize}

\subsection{Status}

\begin{center}
\small
\begin{tabular}{@{}p{0.60\textwidth}p{0.34\textwidth}@{}}
\toprule
Item & Status \\
\midrule
exact ledger identities & compiled \\
finite base & compiled \\
segment and modular exclusions & compiled \\
CRT layer & compiled except the final assembly \\
decisive-window thresholds & unchanged \\
\bottomrule
\end{tabular}
\end{center}
